\scanpage{27}\mspage{26}
\uline{Cont 2} \quad Soit $F$ une figure. \deleted{Alors} Pour tout $X\in F$, \deleted{\unsure{d'une|l'ens.}} \deleted{\unsure{soit}}\unsure{soient|soit}
\[
\partial X^{\deleted{\unsure{\circ|*}}}=\bigcup_{Y\in F_{<X}}Y,\qquad X^\circ=X\setminus\partial X
\]
(strate \uline{ouverte} de $F$ définie par $X$) Alors on a $\boxed{X^\circ\ne\varnothing}$. De plus, si $X,Y\in F$, alors
% slanted
\margin{Les \unsure{ens|ens.} de parties satisfaisant \unsure{ces|aux} conditions seront dits ``admissibles''}
\[
X\cap Y=\bigcup_{\substack{Z\in F\\\unsure{X|Z}\leq\unsure{x|X},Y}}F_{\unsure{Z|X}}
\]
i.e. \deleted{\unsure{tout|tous}} pour \unsure{tt|ttt} \deleted{\unsure{pt}} lieu\textsuperscript{\unsure{$x$}}/ qui appartient à deux strates fermées, \deleted{on a $\exists\,Z\in F$} il existe une \uline{strate} $Z\subset X\cap Y$ \deleted{telle que} \deleted{\unsure{le lieu}} qui contient $x$.

\uline{Corollaire} \deleted{\unsure{Si application|L'application} $F$} \inserted{\deleted{\unsure{par}} L'application,}
\[
\unsure{x|X}\longmapsto X^\circ\qquad F\longrightarrow P(|F|)
\]
\deleted{\unsure{i|in}}induit une bijection de $F$ avec l'ens. des classes d'une relation d'équivalence de $|F|$, i.e pour $X,Y\in F$, \deleted{tels que} \inserted{on a}
\[
\deleted{X\ne Y,\ \text{on a}\quad X^\circ\cap Y^\circ=\varnothing}
\]
\[
X\ne Y\iff X^\circ\cap Y^\circ=\varnothing
\]
(équivaut en fait à \uline{Cont 2})

\deleted{Définition} \inserted{\unsure{Cf.|Cf} \uline{Parties} \uline{saturées} de $|F|$ : p.ex. les $X\in F$}

\deleted{\uline{Définition} Soient $F,\unsure{G|C}$ deux figures. On dit qu'elles sont \uline{compatibles}, si \unsure{pour} les conditions équivalentes a) $|F|\cap|G|$ est saturé pour les relations d'\unsure{équivalence|équivalences} sur $|F|$, $|G|$, et la rel.}
